El desarrollo de sistemas de control por voz en dispositivos de borde ha experimentado un
avance significativo con la evolución de la serie de microcontroladores ESP32. Actualmente,
el ESP32-S3 se posiciona como el estándar para estas aplicaciones debido a su arquitectura de
doble núcleo Xtensa LX7 y, fundamentalmente, por la inclusión de extensiones de
instrucciones de inteligencia artificial que aceleran las operaciones de redes neuronales
(\cite{kumar}). Esta capacidad de procesamiento local permite un cambio de paradigma hacia
la inteligencia en el borde, buscando minimizar la latencia y maximizar la privacidad de los
datos al mantener los procesamientos biométricos de voz dentro del dominio del hardware del
dispositivo.

En cuanto a la captura acústica, se observa una transición desde los micrófonos analógicos
tradicionales hacia los micrófonos MEMS digitales, destacando el modelo INMP441. Estos
dispositivos integran un sensor acústico, un preamplificador y un convertidor analógico a
digital (ADC) de 24 bits, entregando una señal digital directamente a través de la interfaz I2S
(Inter-IC Sound). Este enfoque elimina la vulnerabilidad al ruido inherente a las señales
analógicas. No obstante, investigaciones en entornos educativos todavía emplean soluciones
analógicas como el módulo MAX9814 con control automático de ganancia para simplificar el
diseño de circuitos y reducir costes (\cite{nguyen}).

El procesamiento de voz se divide actualmente en dos grandes arquitecturas: local e
híbrida. El reconocimiento de voz local u offline utiliza marcos de trabajo como ESP-Skainet,
que emplea los motores WakeNet para el monitoreo continuo de palabras de activación y
MultiNet para el reconocimiento de hasta 200 comandos locales (\cite{joseph}).
Complementariamente, plataformas como Edge Impulse facilitan el despliegue de modelos de
aprendizaje automático personalizados (TinyML), entrenados para tareas específicas de
detección de palabras clave con baja latencia. Por otro lado, las arquitecturas híbridas delegan
las tareas computacionalmente intensivas, como el procesamiento de lenguaje natural, a
servicios en la nube como OpenAI Whisper y ChatGPT (GPT-3.5-turbo), permitiendo que
microcontroladores económicos realicen interacciones complejas de voz a texto y texto a voz
(\cite{nguyen}). 

La integración con el Internet de las Cosas (IoT) se apoya predominantemente en el
protocolo MQTT (Transporte de telemetría de colas de mensajes), valorado por su ligereza y
eficiencia en redes con ancho de banda limitado. El uso de MQTT permite que el ESP32 actúe
como un cliente que publica eventos de acceso concedido o se suscribe a comandos remotos,
integrándose con ecosistemas como Home Assistant. Investigaciones emergentes introducen
el protocolo Model Context Protocol (MCP), que actúa como un lenguaje universal para que
modelos de IA interactúen directamente con el hardware, descubriendo capacidades del
dispositivo y ejecutando acciones físicas como el control de cerraduras de solenoide.
En el ámbito de la seguridad, el estado del arte incluye sistemas avanzados de verificación
de identidad mediante embeddings de voz (como ECAPA-TDNN o x-vectors) y mecanismos
de detección de vitalidad (liveness detection) como los sistemas Void y LiveProbe, diseñados
para prevenir ataques de repetición mediante el análisis de la potencia espectral de la voz
humana frente a altavoces. A nivel de hardware, se enfatiza la implementación de Secure
Boot y cifrado de flash para proteger la integridad del firmware y las plantillas biométricas
almacenadas (\cite{jain}).

Finalmente, la implementación práctica enfrenta desafíos críticos en la gestión de
recursos. Las limitaciones de la memoria SRAM (normalmente 512 KB) obligan a optimizar
el uso del heap y emplear técnicas de gestión de tareas mediante FreeRTOS (Sistema
operativo en tiempo real) para manejar de forma concurrente la captura de audio y la
comunicación de red (\cite{nguyen}). Asimismo, la estabilidad eléctrica es un factor
determinante, requiriendo circuitos de protección contra caídas de tensión causadas por la
activación de actuadores de alta corriente como solenoides o ventiladores.